\documentclass[../main-notes.tex]{subfile}

\begin{document}

\justifying

\subsection*{Extra: Basic Trigonometric Functions}

\begin{table}[h]
\centering
\renewcommand{\arraystretch}{1.8}
\begin{tabular}{@{}>{\raggedright\arraybackslash}p{3cm}llp{6cm}@{}}
\toprule
\textbf{Function} & \textbf{Domain} & \textbf{Range} & \textbf{Key Features for Calculus} \\ 
\midrule
\textbf{Sine} & & & \\
\( f(x) = \sin(x) \) & \( (-\infty, \infty) \) & \( [-1, 1] \) & 
\begin{minipage}{6cm}
\begin{itemize}
    \item Continuous everywhere
    \item Derivative: \( \cos(x) \)
    \item Zeros at \( x = n\pi \), \( n \in \mathbb{Z} \)
    \item Period: \( 2\pi \)
\end{itemize}
\end{minipage} \\
\addlinespace
\textbf{Cosine} & & & \\
\( f(x) = \cos(x) \) & \( (-\infty, \infty) \) & \( [-1, 1] \) &
\begin{minipage}{6cm}
\begin{itemize}
    \item Continuous everywhere
    \item Derivative: \( -\sin(x) \)
    \item Zeros at \( x = \frac{\pi}{2} + n\pi \), \( n \in \mathbb{Z} \)
    \item Period: \( 2\pi \)
\end{itemize}
\end{minipage} \\
\addlinespace
\textbf{Tangent} & & & \\
\( f(x) = \tan(x) \) & \( x \neq \frac{\pi}{2} + n\pi \), \( n \in \mathbb{Z} \) & \( (-\infty, \infty) \) &
\begin{minipage}{6cm}
\begin{itemize}
    \item Discontinuous at \( x = \frac{\pi}{2} + n\pi \)
    \item Derivative: \( \sec^2(x) > 0 \) (always increasing on intervals)
    \item Vertical asymptotes at domain exclusions
    \item Period: \( \pi \)
\end{itemize}
\end{minipage} \\
\bottomrule
\end{tabular}
\end{table}

\newpage

\subsection*{General Forms and Transformations}

For functions of the form: \( f(x) = A \sin(Bx + C) + D \) or \( f(x) = A \cos(Bx + C) + D \)

\begin{table}[h]
\centering
\renewcommand{\arraystretch}{1.5}
\begin{tabular}{@{}lll@{}}
\toprule
\textbf{Transformation} & \textbf{Domain} & \textbf{Range} \\ 
\midrule
Amplitude change: \( A\sin(x) \) or \( A\cos(x) \) & \( (-\infty, \infty) \) & \( [-|A|, |A|] \) \\
Vertical shift: \( \sin(x) + D \) or \( \cos(x) + D \) & \( (-\infty, \infty) \) & \( [-1 + D, 1 + D] \) \\
Combined: \( A\sin(x) + D \) or \( A\cos(x) + D \) & \( (-\infty, \infty) \) & \( [-|A| + D, |A| + D] \) \\
Period change: \( \sin(Bx) \) or \( \cos(Bx) \) (\( B > 0 \)) & \( (-\infty, \infty) \) & \( [-1, 1] \) \\
Phase shift: \( \sin(x + C) \) or \( \cos(x + C) \) & \( (-\infty, \infty) \) & \( [-1, 1] \) \\
\bottomrule
\end{tabular}
\end{table}

\subsection*{General Tangent Function}

For \( f(x) = A \tan(Bx + C) + D \):
\begin{itemize}
    \item \textbf{Domain}: Solve \( Bx + C \neq \frac{\pi}{2} + n\pi \), \( n \in \mathbb{Z} \)
    \[
    x \neq \frac{1}{B}\left(\frac{\pi}{2} + n\pi - C\right)
    \]
    \item \textbf{Range}: \( (-\infty, \infty) \) (unchanged by vertical shifts and stretches)
    \item \textbf{Period}: \( \frac{\pi}{|B|} \)
\end{itemize}

\subsection*{Calculus Applications and Important Limits}

\begin{table}[h]
\centering
\renewcommand{\arraystretch}{1.5}
\begin{tabular}{@{}lp{10cm}@{}}
\toprule
\textbf{Limit} & \textbf{Significance in Calculus I} \\ 
\midrule
\( \displaystyle \lim_{x \to 0} \frac{\sin(x)}{x} = 1 \) & 
Fundamental limit for derivatives of trig functions, used in squeeze theorem examples \\
\addlinespace
\( \displaystyle \lim_{x \to 0} \frac{1 - \cos(x)}{x} = 0 \) &
Used in derivative of cosine function derivation \\
\addlinespace
\( \displaystyle \lim_{x \to \frac{\pi}{2}^-} \tan(x) = \infty \) &
Example of infinite limit and vertical asymptote behavior \\
\addlinespace
\( \displaystyle \lim_{x \to \frac{\pi}{2}^+} \tan(x) = -\infty \) &
Important for understanding one-sided limits at discontinuities \\
\bottomrule
\end{tabular}
\end{table}

\subsection*{Derivatives and Continuity}

\begin{itemize}
    \item \textbf{Sine and Cosine}:
    \begin{itemize}
        \item Continuous everywhere, differentiable everywhere
        \item \( \frac{d}{dx}[\sin(x)] = \cos(x) \)
        \item \( \frac{d}{dx}[\cos(x)] = -\sin(x) \)
    \end{itemize}
    
    \item \textbf{Tangent}:
    \begin{itemize}
        \item Continuous and differentiable on its domain
        \item \( \frac{d}{dx}[\tan(x)] = \sec^2(x) \)
        \item Discontinuous at \( x = \frac{\pi}{2} + n\pi \), vertical asymptotes at these points
    \end{itemize}
    
    \item \textbf{Chain Rule Applications}:
    \[
    \frac{d}{dx}[\sin(u(x))] = \cos(u(x)) \cdot u'(x)
    \]
    \[
    \frac{d}{dx}[\tan(u(x))] = \sec^2(u(x)) \cdot u'(x)
    \]
\end{itemize}

\subsection*{Important Properties for Problem Solving}

\begin{enumerate}
    \item \textbf{Sine and Cosine}:
    \begin{itemize}
        \item Bounded between -1 and 1, useful for squeeze theorem
        \item Periodic nature affects behavior of derivatives and integrals
        \item \(\sin^2(x) + \cos^2(x) = 1\) (Pythagorean identity, useful in simplifications)
    \end{itemize}
    
    \item \textbf{Tangent}:
    \begin{itemize}
        \item Always increasing on each interval of its domain
        \item Related to sine and cosine: \( \tan(x) = \frac{\sin(x)}{\cos(x)} \)
        \item Derivative is always positive: \( \sec^2(x) > 0 \)
    \end{itemize}
    
    \item \textbf{Domain Considerations in Calculus}:
    \begin{itemize}
        \item When differentiating composite functions with trig functions inside, ensure inner function is in domain
        \item For \(\tan(f(x))\), require \( f(x) \neq \frac{\pi}{2} + n\pi \)
        \item When finding critical points, consider full period (not just \([0, 2\pi]\))
    \end{itemize}
\end{enumerate}

\subsection*{Graphical Behavior}

\begin{itemize}
    \item \textbf{Sine and Cosine}: Smooth, wavelike graphs (sinusoids)
    \begin{itemize}
        \item Sine: Starts at 0, maximum at \(\pi/2\), zero at \(\pi\), minimum at \(3\pi/2\)
        \item Cosine: Starts at 1, zero at \(\pi/2\), minimum at \(\pi\), zero at \(3\pi/2\)
    \end{itemize}
    
    \item \textbf{Tangent}: Discontinuous graph with vertical asymptotes
    \begin{itemize}
        \item Period of \(\pi\) (vs \(2\pi\) for sine/cosine)
        \item Passes through origin, repeats every \(\pi\) units
        \item Approaches \(\pm\infty\) near vertical asymptotes
    \end{itemize}
\end{itemize}


\end{document}
